\clearpage
\section{证据索引与文献}
\subsection{内部证据：读取范围、版本和资格}
固定科学仓库为 \url{https://github.com/migodam/b1-probabilistic-maxwell-imaging}，
版本为 \texttt{d1bdc8aceae4cd6a843f873b61b019d5eee3cd1d}。
本轮读取了上传 ZIP 中用户指定的 R6 主报告、理论反例、各阶段报告、GO/NO-GO 和 Handoff，并读取了下列固定版本历史原报告。R3 和 R5T 的完整独立 Pro 报告没有取得：涉及其主张的部分采用本次用户提供的明确陈述，以及 R5/R6 对它们的引用；不将这种二级追溯写成完整原文复核。
\begin{longtable}{p{.08\linewidth}p{.84\linewidth}}
\toprule 编号 & 来源与用途\\\midrule\endhead
H1 & \texttt{research\_lxz\_20261016}：L/X/Z 非线性探索、噪声相容性、proxy 比较。\\
H2 & \texttt{research\_f\_intrinsic\_transport}：decoder 二阶贡献、Maxwell 二阶向量证据缺口、F2 站稳条件导致的未执行。\\
H3 & \texttt{research\_f3\_quotient\_exploration}：公共权重尺度商、非均匀权重比先验、局部可行端点与路径限制。\\
H4 & \texttt{a17x\_f3\_competition\_audit}：七维参数与高秩材料增量、端口截断代价、A17 适用范围。\\
H5 & \texttt{p3l\_local\_posterior\_imaging}：动态弱方向、局部 posterior provider、退化参考与材料风险。\\
H6 & \texttt{p3l\_r2\_risk\_posterior\_exploration}：独立材料风险、Maxwell 特征对照、真实材料端点、非对合证据。\\
H7 & \texttt{p3l\_r4\_conditional\_transport}：Full-66 相对概率指标、绝对资格失败、独立材料风险和非完整路径。\\
H8 & \texttt{p3l\_r5\_parallel\_ab}：分层概率与蒸馏阶段报告。R5-B 后续 512 对象确认另见 R6-D 引用的 prior evidence 和审计，不以本阶段尚未运行的计划代替确认结果。\\
R6-A & \texttt{THEORY\_AND\_COUNTEREXAMPLES.md}；\newline 
\texttt{STAGE\_A\_STRATIFIED\_PROBABILITY.md}；\newline 
\texttt{STAGE\_A\_MAXWELL\_COMPENSATION.md}；\newline 
\texttt{STAGE\_A\_TASK\_INTEGRATION.md}。\\
R6-B & \texttt{STAGE\_B\_FUNCTIONAL\_DISTILLATION.md}；\newline 
\texttt{STAGE\_B\_SEED\_VARIANCE\_AND\_ACTION\_RISK.md}；\newline 
\texttt{STAGE\_B\_MAXWELL\_CONTEXT\_ABLATION.md}。\\
R6-C & \texttt{STAGE\_C\_PHYSICS\_ANCHOR\_COST.md}；\newline 
\texttt{STAGE\_C\_RISK\_AWARE\_STOPPING.md}；\newline 
\texttt{STAGE\_D\_LOCKED\_GENERALIZATION.md}；\newline 同窗口复测、Direct Probit 和 Full-q 强基线补充审计。\\
R6-D & \texttt{RESEARCH\_REPORT\_ZH.md}；\newline 
\texttt{GO\_NO\_GO.md}；\newline 
\texttt{GPT\_PRO\_HANDOFF/START\_HERE.md}；\newline 
\texttt{PRIOR\_EVIDENCE\_AND\_FAILURES.md}；\newline 
machine summaries 与审计。六种 locked 方法、两噪声档的 12 组已有 MSE 数组经过独立重算；没有重跑物理。\\
A17 & \texttt{migodam/a17-maxwell-certified-exchange} 中
\texttt{closure\_final/manuscript/CORE\_THEORY.tex} 及
\texttt{CONDITIONAL\_EXCHANGE\_THEORY.tex}。
读取 blob 为
\texttt{c31c5c6790820b153d2b6a738df94fda0c12165d} 和
\texttt{ee20bd2e00951a1c3adbaa9cfc44cfcdb2db4b62}。
用于反馈、current 闭合与双侧任务传递核查，不视为已发表外部先例。\\
\bottomrule
\end{longtable}

\subsection{原始文献及其对本研究的限制}
以下采用作者原文、期刊页面、大学仓储或 arXiv 原始稿。2026 年新预印本明确标出。检索截止日不等于证明不存在其它先例。

\begin{enumerate}[label=\textbf{L\arabic*.},leftmargin=3em,itemsep=.7em]
\item X. Chen, “Subspace-Based Optimization Method for Solving Inverse-Scattering Problems,” \emph{IEEE TGRS}, 48(1), 42--49, 2010.
\url{https://doi.org/10.1109/TGRS.2009.2025122}。
作者列表：\url{https://www.ece.nus.edu.sg/stfpage/elechenx/Publications.htm}。
\textbf{边界：} 数据确定与歧义 current 分量、子空间约束及物理闭合已有先例。

\item Y. Zhong and X. Chen, “Twofold Subspace-Based Optimization Method for Solving Inverse Scattering Problems,” \emph{Inverse Problems}, 25, 085003, 2009.
\url{https://doi.org/10.1088/0266-5611/25/8/085003}。
\textbf{边界：} TSOM 已用双重子空间构造和内部物理约束进一步处理 current 歧义，但不是材料后验分支归一化。

\item S. Arhab, D. Anagnostou and M. Joelson, “High-Order Functional Derivatives of the Diffracted Field According to the Permittivity-Contrast Function,” \emph{Wave Motion}, 83, 67--79, 2018.
\url{https://doi.org/10.1016/j.wavemoti.2018.07.008}；
\url{https://arxiv.org/abs/1803.00002}。
\textbf{边界：} 高阶散射导数及递推已有。增量须落在补偿后的首个可观测系数、有限误差窗口和采集冲突上。

\item Y. Liu et al., “SOM-Net: Unrolling the Subspace-Based Optimization for Solving Full-Wave Inverse Scattering Problems,” \emph{IEEE TGRS}, 60, article 2007715, 2022.
\url{https://doi.org/10.1109/TGRS.2022.3217162}；
\url{https://arxiv.org/html/2209.03567v2}。
\textbf{边界：} current/contrast 物理网络已有；R7 要证明材料任务信息或相容材料概率体积的额外价值。

\item T. Cui et al., “Likelihood-Informed Dimension Reduction for Nonlinear Inverse Problems,” \emph{Inverse Problems}, 30, 114015, 2014.
\url{https://arxiv.org/abs/1403.4680}。
\textbf{边界：} posterior 信息几何低维化已有成熟方法。

\item O. Zahm et al., “Certified Dimension Reduction in Nonlinear Bayesian Inverse Problems,” \emph{Mathematics of Computation}, 91, 1789--1835, 2022.
\url{https://doi.org/10.1090/mcom/3737}；
\url{https://arxiv.org/abs/1807.03712}。
\textbf{边界：} 导数驱动概率降维及误差控制已有。须说明原子、重叠材料像或任务误差的新增困难。

\item A. Spantini, T. Cui, K. Willcox, L. Tenorio and Y. Marzouk, “Goal-Oriented Optimal Approximations of Bayesian Linear Inverse Problems,” \emph{SIAM J. Scientific Computing}, 39(5), S167--S196, 2017.
\url{https://doi.org/10.1137/16M1082123}；
\url{https://arxiv.org/abs/1607.01881}。
\textbf{边界：} 目标协方差与 Bayes-risk 最优低秩逼近是直接先例。

\item J. T. Chang and D. Pollard, “Conditioning as Disintegration,” \emph{Statistica Neerlandica}, 51(3), 287--317, 1997.
\url{https://doi.org/10.1111/1467-9574.00056}；
\url{https://www.cs.tufts.edu/comp/150FP/archive/david-pollard/ConditioningAsDisintegration.pdf}。
\textbf{边界：} 条件核与一般测度分解是成熟基础；命题 1 不是新的 disintegration 定理。

\item P. Diaconis, S. Holmes and M. Shahshahani, “Sampling From a Manifold,” \emph{IMS Collections}, 10, 102--125, 2013.
\url{https://doi.org/10.1214/12-IMSCOLL1006}；
\url{https://arxiv.org/abs/1206.6913}。
\textbf{边界：} level-set 条件测度和 coarea 权重已有；均匀轨迹点不等于目标分布。

\item K. Kilgore, S. Moskow and J. C. Schotland, “Convergence of the Born and Inverse Born Series for Electromagnetic Scattering,” \emph{Applicable Analysis}, 96(10), 1737--1748, 2017.
\url{https://doi.org/10.1080/00036811.2017.1292349}。
\textbf{边界：} Maxwell Born 级数和收敛理论已有，非零背景尾界仍须可验证常数。

\item J. M. Bardsley, A. Solonen, H. Haario and M. Laine, “Randomize-Then-Optimize: A Method for Sampling from Posterior Distributions in Nonlinear Inverse Problems,” \emph{SIAM J. Scientific Computing}, 36(4), A1895--A1910, 2014.
\url{https://doi.org/10.1137/140964023}；
\url{https://dspace.mit.edu/entities/publication/ead8cb08-a5cb-4320-8260-135828ed5745}。
\textbf{边界：} 优化建议与概率校正已有；候选多样性不能替代权重。

\item T. O'Leary-Roseberry, P. Chen, U. Villa and O. Ghattas, “Derivative-Informed Neural Operator: An Efficient Framework for High-Dimensional Parametric Derivative Learning,” \emph{J. Computational Physics}, 496, 112555, 2024.
\url{https://doi.org/10.1016/j.jcp.2023.112555}；
\url{https://arxiv.org/abs/2206.10745}。
\textbf{边界：} 导数信息、低维编码和算子学习已有；比较须计入离线物理费。

\item L. Cao, T. O'Leary-Roseberry and O. Ghattas, “Derivative-Informed Neural Operator Acceleration of Geometric MCMC for Infinite-Dimensional Bayesian Inverse Problems,” \emph{JMLR}, 26(78), 1--68, 2025.
\url{https://jmlr.org/papers/v26/24-0745.html}。
\textbf{边界：} DINO 校正采样是直接竞争路线。

\item L. Cao et al., “LazyDINO: Fast, Scalable, and Efficiently Amortized Bayesian Inversion via Structure-Exploiting and Surrogate-Driven Measure Transport,” \emph{JMLR}, 27(24), 1--71, 2026.
\url{https://jmlr.org/papers/v27/25-0858.html}。
\textbf{边界：} derivative learning、lazy transport 与校正推断已有系统框架；不能只提出 derivative-aware Flow。

\item S. Watanabe, “A Widely Applicable Bayesian Information Criterion,” \emph{JMLR}, 14(27), 867--897, 2013.
\url{https://www.jmlr.org/papers/v14/watanabe13a.html}。
\textbf{边界：} 奇异积分的幂次与对数修正已有。一条材料弧不决定高维 posterior 通用指数。

\item J. Xu, Y. Wu, D. Zeng, J. Paisley and Q. Zhao, “The Right Measure for Physics-Constrained Generation: A Co-Area Correction for Posterior-Consistent PDE Inverse Problems,” arXiv:2606.04804v3, 2026 年 6 月 7 日，\textbf{预印本}。
\url{https://arxiv.org/html/2606.04804v3}。
\textbf{边界：} 其正则约束、连续正先验和满行秩假设限制“coarea sampling 是新方法”的主张；原子和秩变接缝仍需另处理。

\item A. Dytso, H. V. Poor and S. Shamai, “Conditional Mean Estimation in Gaussian Noise: A Meta Derivative Identity With Applications,” \emph{IEEE Transactions on Information Theory}, 69(3), 1883--1898, 2023.
\url{https://doi.org/10.1109/TIT.2022.3216012}；
\url{https://arxiv.org/abs/2104.01883}。
\textbf{边界：} Gaussian conditional score/covariance identity 已知；命题 4 的一般算子不是新的 Maxwell 定理。

\item S. Zhong et al., “Goal-Oriented Bayesian Optimal Experimental Design for Nonlinear Models Using Markov Chain Monte Carlo,” arXiv:2403.18072v2, 2025 年 2 月 2 日作者稿。
\url{https://arxiv.org/html/2403.18072v2}；
\url{https://doi.org/10.1137/24M1649344}。
\textbf{边界：} 非线性任务导向采集已有。这里不依赖未完整核对的正式卷页元数据。

\item J. D. Jakeman, R. White, B. van Bloemen Waanders, D. P. Kouri and A. Alexanderian, “Risk-Aware Goal-Oriented Bayesian Optimal Experimental Design,” arXiv:2609.08023v1, 2026 年 9 月 7 日，\textbf{预印本}。
\url{https://arxiv.org/html/2609.08023v1}。
\textbf{边界：} 风险导向设计不是空白领域；新增指标须比较真实任务 Bayes-risk。

\item J. Wen, R. Ahmad and P. Schniter, “Task-Driven Uncertainty Quantification in Inverse Problems via Conformal Prediction,” \emph{ECCV}, 182--199, 2024.
\url{https://doi.org/10.1007/978-3-031-73027-6_11}；
\url{https://arxiv.org/html/2405.18527v2}。
\textbf{边界：} 任务 UQ 可不依赖完整 posterior；边际覆盖不等于逐观测 posterior 正确。

\item P. Kerfriden, A. Kundu and S. Claus, “Adaptivity in Bayesian Inverse Finite Element Problems: Learning and Simultaneous Control of Discretisation and Sampling Errors,” \emph{Materials}, 12(4), 642, 2019.
\url{https://doi.org/10.3390/ma12040642}；
\url{https://orca.cardiff.ac.uk/id/eprint/119861/1/materials-423575.pdf}。
\textbf{边界：} 离散与采样误差联控已有；不能只将误差预算拆分就认领原创性。

\item S. L. Cotter, M. Dashti and A. M. Stuart, “Approximation of Bayesian Inverse Problems for PDEs,” \emph{SIAM J. Numerical Analysis}, 48(1), 322--345, 2010.
\url{https://doi.org/10.1137/090770734}；
\url{https://arxiv.org/html/0909.2126v1}。
\textbf{边界：} forward 逼近到 posterior 与函数量的误差传播已有。

\item J. A. Christen and C. Fox, “Markov Chain Monte Carlo Using an Approximation,” \emph{J. Computational and Graphical Statistics}, 14(4), 795--810, 2005.
\url{https://doi.org/10.1198/106186005X76983}。
\textbf{边界：} delayed acceptance 的廉价筛选与精确校正已有。

\item Z. Han, Y. Liang and H. Zhao, “Iterative Truncated SVD for Inverse Scattering Problems,” arXiv:2610.25726v1, 2026 年 10 月 19 日，\textbf{预印本}。
\url{https://arxiv.org/html/2610.25726v1}。
\textbf{边界：} iTSVD 是动态截断谱路线的直接新近竞争；不能将其性能主张当作本报告独立验证的事实。

\item X. Ye and X. Chen, “Subspace-Based Distorted-Born Iterative Method for Solving Inverse Scattering Problems,” \emph{IEEE TAP}, 65(12), 7224--7232, 2017.
\url{https://doi.org/10.1109/TAP.2017.2766658}；
\url{https://pure.bit.edu.cn/en/publications/subspace-based-distorted-born-iterative-method-for-solving-invers/}。
\textbf{边界：} 子空间与非线性背景更新结合已有成熟 TAP 先例。

\item T. Butler, D. Estep, S. Tavener, T. Wildey, C. Dawson and L. Graham, “Solving Stochastic Inverse Problems Using Sigma-Algebras on Contour Maps,” arXiv:1407.3851v1, 2014，作者预印本。
\url{https://arxiv.org/abs/1407.3851}。
\textbf{边界：} contour/quotient sigma-algebra 组织逆概率已有；其测度设定不应自动等同于指定 prior 加 likelihood。
\end{enumerate}

\subsection{可追溯性说明}
Disintegration、coarea、隐函数约化、条件期望投影、Bayes-risk 恒等式、adjoint 残差和分层最优分配均有成熟来源。独立写出命题是为了证明本问题的适用边界，不是认领全部公式为新定理。

最值得继续验证的组合是：同一非均匀材料背景与同一可行补偿空间上，以真实 Maxwell current 递推得到首个可用残差系数及有限幅度区间；再用共享补偿冲突识别新采集的物理价值。是否足够原创，取决于与高阶散射、profile/Schur information 和任务设计的比较。

交付包包含可编辑 \LaTeX 源、解析图代码和独立 Codex 任务书。解析图只解释公式，不能替代尚未执行的 Maxwell 实验。

